\chapter{系统测试与分析}

本章将对 Chimera 渲染引擎进行全面的系统测试。测试不仅涵盖混合渲染管线的视觉质量分析，还将通过定量的性能指标评估系统的实时性。此外，本章通过对比实验验证了 SVGF 降噪算法与 TAA 抗锯齿模块在提升画面质量方面的关键作用。

\section{实验环境与方案设计}

为了确保测试数据的客观性与可重复性，本系统的所有渲染效果捕获与定性分析均在统一的软硬件环境下进行。

系统的基础硬件配置为：处理器采用 AMD Ryzen 9 9955HX 16-Core Processor，图形处理器为支持硬件级光线追踪的 NVIDIA GeForce RTX 5070 Ti，配备 32GB 系统主存。

软件环境基于 Windows 11 操作系统，图形 API 采用 Vulkan 1.3 SDK，着色器语言采用 GLSL（编译为 SPIR-V 字节码）。系统的渲染结果截取与画面层级分析主要依赖于 RenderDoc 图形调试工具。

考虑到现代图形学对高分辨率显示的需求，本系统的所有视觉测试基准分辨率统一设定为 2560$\times$1440 (2K 分辨率)。测试场景选取了图形学界经典的两个基准模型：
\begin{enumerate}
    \item \textbf{Sponza 宫殿场景}：包含复杂的室内几何结构与多种 PBR 纹理（共计约 26 万个三角形），用于极限测试光追阴影的遮挡计算与 SVGF 在复杂纹理下的降噪表现。
    \item \textbf{多材质反射测试场景（Wuson \& Spheres）}：包含具有不同粗糙度（Roughness）与金属度（Metallic）的 PBR 材质物体，用于验证 RT Pipeline 镜面反射的物理正确性。
\end{enumerate}

\section{混合渲染质量评估}

本节对混合渲染管线的视觉表现进行分析。

\subsection{G-Buffer 中间数据准确性验证}
作为混合渲染的数据基底，G-Buffer 提取的准确性直接决定了后续光照与降噪的正确性。如第五章图 \ref{fig:gbuffer_outputs} 所示，反照率（Albedo）缓冲精确捕捉了材质的基础颜色；世界空间法线（Normal）缓冲呈现出平滑的几何法向过渡；运动矢量（Motion Vector）缓冲则准确记录了摄像机平移时像素在屏幕空间的二维物理位移。这些纯净的几何与材质信号为后续的光追阶段奠定了坚实的基础。

\subsection{光线追踪特效分析}
在传统的光栅化管线中，阴影通常依赖级联阴影贴图（CSM），而反射则依赖屏幕空间反射（SSR）。本系统采用硬件光线追踪替代了这些经验算法：
\begin{itemize}
    \item \textbf{硬软阴影过渡}：通过 Ray Query 发射的可见性射线，系统彻底消除了 CSM 常见的锯齿与漏光现象。结合面法线偏导数的偏移算法，模型表面未出现任何自相交伪影。
    \item \textbf{物理精确的镜面反射}：在多材质球体场景中，基于 RT Pipeline 追踪的反射射线不仅能够完美倒影出屏幕空间外（Off-screen）的物体，彻底解决了 SSR 的断裂伪影，且根据 GGX 微面元分布函数，不同粗糙度的球体呈现出极其真实的物理高光衰减。
\end{itemize}

\subsection{混合渲染与传统光栅化对比}
通过将传统的 Forward Rasterization（仅包含基础冯氏光照与经验贴图）与本文实现的混合渲染管线进行对比，可以看出：
\begin{enumerate}
    \item \textbf{阴影质量}：混合渲染呈现了基于物理的接触遮挡（Contact Hardening）效果，而传统光栅化在几何重叠处阴影缺失严重。
    \item \textbf{间接光照与环境遮蔽}：混合渲染通过光追实现了更为自然的地面环境遮蔽（AO），消除了物体“漂浮”感。
    \item \textbf{材质表现力}：混合渲染下的金属材质能够实时反射周围环境，增强了画面的沉浸感。
\end{enumerate}

\section{SVGF 降噪与 TAA 抗锯齿性能评估}

由于实时性能的严格限制，本系统光追阶段的初始采样率仅为 1 SPP。

\subsection{SVGF 降噪稳定性分析}
在禁用降噪模块时，1 SPP 原始输出画面中存在强烈的蒙特卡洛高频噪声。如第五章图 \ref{fig:svgf_result} 所示，经过 SVGF 模块处理后：
\begin{enumerate}
    \item \textbf{反照率解耦的有效性}：由于光照信号在进入 SVGF 前已被剥离了 Albedo，空域滤波并未将材质纹理误判为噪声。
    \item \textbf{深度线性化的有效性}：在场景深处，由于采用了视图空间的线性深度进行历史一致性校验，系统并未产生错误的遮挡剔除，画面中未观察到明显的历史残影（Ghosting）。
\end{enumerate}

\subsection{TAA 抗锯齿有效性验证}
通过在静止相机下对高频几何模型进行对比观察。如第五章图 \ref{fig:taa_result} 所示，TAA 的引入不仅消除了几何边缘的阶梯状走样（Aliasing），更重要的是通过时域累积增强了画面的稳定性。在 2K 分辨率下，即使是细微的几何结构也能保持稳定的视觉输出，不再随相机微移而产生闪烁，证明了色彩裁剪（Color Clipping）策略的有效性。

\section{系统集成界面与交互功能展示}

除了底层渲染算法的实现，本系统还开发了一套基于 ImGui 增强的实时编辑器界面。

编辑器界面的核心功能包括：
\begin{itemize}
    \item \textbf{动态管线配置}：用户可以在运行时无缝切换 Forward、Deferred 或 Hybrid 渲染路径。
    \item \textbf{参数在线调试}：支持对 SVGF 滤波半径、TAA 混合因子、光追递归深度等关键算法参数进行滑块控制。
    \item \textbf{资源实时监控}：集成的资源管理器可实时反馈显存占用、纹理加载状态以及每帧 GPU 耗时，确保了系统的高效运行。
\end{itemize}

\section{本章小结}

本章通过定性和定量的对比分析，评估了本课题实现的实时光追混合渲染引擎。实验结果表明，系统成功还原了物理正确的阴影软硬过渡与复杂的多次镜面反射。经过优化的 SVGF 与 TAA 算法，在不增加额外光线发射数量的前提下，显著提升了极端欠采样条件下的画面信噪比与几何稳定性。本章的测试结果有效验证了前文理论推导与工程设计方案的正确性。
